\documentclass{article}

%----------------------------------------
% Packages
%----------------------------------------
\usepackage[left=1in, right=1in, top=1in, bottom=1in]{geometry}
\usepackage{graphicx}
\usepackage{amsmath,amsbsy,amssymb,amsfonts,amsthm}
\usepackage{nicefrac}
\usepackage{mathtools}
\usepackage{color}
\usepackage{xspace} % Correct macro spacing
\usepackage[numbers]{natbib} % For citations
\usepackage{times}
\usepackage{graphicx,subfigure}
%\usepackage[small,bf]{caption}
\usepackage{algorithm,algorithmic} 
\usepackage{hyperref}
\hypersetup{pdftitle={Stability and generalization: lecture notes, IFT 6169 (formerly IFT 6085), Université de Montréal},
            pdfauthor={Ioannis Mitliagkas (Γιάννης Μητλιάγκας), instructor, and student scribes}}

\usepackage{xcolor}
\usepackage{shadethm}

\usepackage{fancyhdr}
\pagestyle{fancy}
\lhead{This is my name}
\rhead{this is page \thepage}

\usepackage{fancyhdr}
\pagestyle{fancy}
\lhead{IFT 6085 - Theoretical principles for deep learning}
\rhead{Lecture 11: February 20, 2019}


\newshadetheorem{thm}{Theorem}
\newshadetheorem{defn}[thm]{Definition}
\newshadetheorem{assm}[thm]{Assumption}
\newshadetheorem{prop}[thm]{Property}
\newshadetheorem{eg}[thm]{Example}
% Define a Lemma env.
\newtheorem{theorem}{Theorem}[section]
\newtheorem{lemma}[theorem]{Lemma}
\usepackage{xcolor,colortbl}


\definecolor{shadethmcolor}{HTML}{F0F0F0}
%\definecolor{shadethmcolor}{HTML}{EDEDED}
%\definecolor{shadethmcolor}{HTML}{EDF8FF}
%\definecolor{shaderulecolor}{HTML}{EDF8FF}
%\definecolor{shaderulecolor}{HTML}{45CFFF}
\setlength{\shadeboxrule}{.4pt}


\setlength\parindent{0pt}

% Packages hyperref and algorithmic misbehave sometimes.  We can fix
% this with the following command.
\providecommand{\theHalgorithm}{\arabic{algorithm}}
\newcommand{\amy}[1]{\textbf{\textcolor{blue}{#1}}}
\newcommand{\liam}[1]{\textbf{\textcolor{red}{#1}}}

%----------------------------------------
% Standard macros
%----------------------------------------


%----------------------------------------
% Project-specific macros
%----------------------------------------

%----------------------------------------
% Header
%----------------------------------------
\title{IFT 6085 - Lecture 11 \\ 
(Applications of Stability) }
\date{}

%----------------------------------------
% Document
%----------------------------------------
\begin{document} 

%----------------------------------------
% Abstract
%----------------------------------------
\maketitle

\textbf{Scribes:} 

\textbf{Winter 2018}: Amy Zhang, William Fedus
\hfill
\textbf{Instructor:} {\hypersetup{pdfborder={0 0 0}}\href{https://mitliagkas.github.io/}{Ioannis Mitliagkas}}

\textbf{Winter 2019}: Bhargav Kanuparthi, Jonathan Pilault


%----------------------------------------
% Body
%----------------------------------------

\newcommand{\infgc}{\inf_{g \in \mathcal{C}}}
\newcommand{\supgc}{\sup_{g \in \mathcal{C}}}

\newcommand{\Prob}{\mathbb{P}}
\newcommand{\E}{\mathbb{E}}
\newcommand{\reals}{\mathbb{R}}

\section*{Summary}
\textbf{Sufficient condition: Given enough samples we can achieve a good enough generalization.  However, typically in deep learning, we never have large enough data sets to get non-vacuous or meaningful bounds.}

\begin{table}[h!]
\centering
\begin{tabular}{|c  | c|} 
 \hline
 \rowcolor{shadethmcolor}
 \textbf{Last Time} & \textbf{Today} \\ [0.5ex] 
 \hline
 PAC Bounds &  Stability \\
 Occam Bounds & PAC Bayes \\
 PAC Bayes Bounds & (Practical) Generalization \\
 Stability Bounds &  \\
\hline
\end{tabular}
\end{table}

How can we go from PAC Bayes to a non-vacuous generalization bound?

\vspace{0.3cm}
By sacrificing some data as part of a dedicated test set, we can measure test set generalization and achieve a tighter bound than the weak population bounds.  See \textit{Tutorial on Practical Prediction Theory for Classification} \cite{Langford:2005:TPP:1046920.1058111} for a comprehensive examination.


\section*{Stability}
\begin{defn}[Uniformly $\beta$-stable algorithm]
\[
h_s = \mathcal{A}(S), h_s\in\mathcal{H}
\]
Algorithm $\mathcal{A}$ is stable if $\forall (S, z), \forall i=\{1,...,n\}$
\[
\sup_{z'\in\mathcal{Z}}|l(h_s, z') - l(h_{s_{i,z}}, z')| \leq \beta
\]

\noindent where $S$ is the data set, $z$ is an evaluation sample and $S_{i,z}$ refers to replacing the $i^{th}$ element in $S$ with $z$.  
\end{defn}



\begin{thm}
Consider a $\beta$-uniformly stable algorithm $\mathcal{A}$ with respect to a loss function $\ell : \mathcal{Y} \times \mathcal{Y} \rightarrow [0, M]$ and a hypothesis $h_s$ with $|S|=n$. The following bound holds with a probability of $1-\delta$.

\[
R[h_s]\leq\hat{R}_s[h_s] + \beta  + (\beta n + \frac{M}{2})\sqrt{\frac{2\ln{2/\delta}}{n}}
\]
\end{thm}

The term $(\beta n + \frac{M}{2})\sqrt{\frac{2\ln{2/\delta}}{n}}$ is $O(\beta\sqrt{n})$. Informally, an algorithm is stable if $\beta=O(\frac{1}{n})$. If stability is $O(\frac{1}{\sqrt{n}})$, this term is $O(1)$ and we can no longer show decrease in generalization gap with with increase in $n$.

\section*{Empirical Risk Minimization + Regularization is Stable}
Notation: 
$$\hat{R}_S(w) \triangleq \hat{R}_S(h_w)$$ where $h_w$ is a model parameterized by weights $w$.
\begin{align*}
    l(h,z)&\equiv l(h(x), y) \\
    l(h_w, z) &\equiv l(w,z)
\end{align*}

\begin{thm}[ERM with regularization is $\beta$-stable]
Under the assumption that $\hat{R}_S(w)$ is convex and $l(\cdot|z)$ is L-Lipschitz $\forall z$, Empirical Risk Minimization and Regularization is $\beta$ uniformly stable where

\[
\beta = \frac{4L^2}{\lambda n}
\]
\end{thm}
\begin{proof}
The objective function to be optimized can be written as
$$f_S(w)=\hat{R}_S(w) + \frac{\lambda}{2}||w||_2^2$$
Consider weights $u,v$ for two different models.
$$f_S(v) - f_S(u) = [\hat{R}_S(v) + \frac{\lambda}{2}||v||_2^2] - [\hat{R}_S(u) + \frac{\lambda}{2}||u||_2^2]$$
We perturb the dataset by replacing the data point at $i$ with $z_i'$. Now we get:
\begin{align*}
    f_S(v) - f_S(u) &= \hat{R}_{S_{i,z_i'}}(v) + \frac{\lambda}{2}||v||_2^2 - (\hat{R}_{S_{i,z_i'}}(u) + \frac{\lambda}{2}||u||^2_2) + \frac{l(v, z_i) - l(v, z_i')}{n} - \frac{l(u, z_i) - l(u,z_i')}{n} \\
    &= f_{S_{i,z_i'}}(v) - f_{S_{i,z_i'}}(u) + \frac{l(v, z_i) - l(v, z_i')}{n} - \frac{l(u, z_i) - l(u,z_i')}{n}
\end{align*}
Now we substitute $v=\mathcal{A}(S_{i,z_i'})$ and $u=\mathcal{A}(S)$.
\begin{align*}
    f_S(\mathcal{A}(S_{i,z_i'})) - f_S(\mathcal{A}(S)) 
    =& f_{S_{i,z_i'}}(\mathcal{A}(S_{i,z_i'})) - f_{S_{i,z_i'}}(\mathcal{A}(S)) \\ &+ \frac{l(\mathcal{A}(S_{i,z_i'}), z_i) - l(\mathcal{A}(S_{i,z_i'}), z_i')}{n} - \frac{l(\mathcal{A}(S), z_i) - l(\mathcal{A}(S),z_i')}{n}
\end{align*}

Because
\begin{align*}
f_{S_{i,z_i'}}(\mathcal{A}(S_{i, z_i'})) = \min_w f_{S_{i,z_i'}}(w) \\
\implies \forall w f_{S_{i,z_i'}}(w) \geq f(S_{i,z_i'})(\mathcal{A}(S_{i, z_i'}))
\end{align*}

\begin{assm}
    $l(\cdot|z)$ is L-Lipschitz.
\end{assm}

\begin{align}
f_S(\mathcal{A}(S_{i,z_i'})) - f_S(\mathcal{A}(S)) &\leq \frac{l(\mathcal{A}(S^{i,z_i'}), z_i) - l(\mathcal{A}(S), z_i)}{n} - \frac{l(\mathcal{A}(S^{i,z_i'}), z_i') - l(\mathcal{A}(S),z_i')}{n} \nonumber\\
&\leq 2\frac{L}{n}||\mathcal{A}(S)-\mathcal{A}(S_{i,z_i'})||_2 \label{eq:one}
\end{align}

\begin{assm}
    $\hat{R}_S(w)$  is cvx.
\end{assm}
Which gives us $f_S(w)$ is $\lambda$-str cvx.
Now we perform a Taylor expansion:
\begin{align}
    f_S(\mathcal{A}(S_{i,z_i'})) - f_S(\mathcal{A}(S)) \geq \frac{\lambda}{2}||\mathcal{A}(S_{i,z_i'}) - \mathcal{A}(S)||_2^2
\label{eq:two}
\end{align}
Since $\mathcal{A}(S)$ is the minimizer of $f_s$ and $\lambda$-str cvx the first term disappears.

From \ref{eq:one} and \ref{eq:two} we get:
\begin{align}
    ||\mathcal{A}(S) - \mathcal{A}(S_{i,z_i'})|| \leq \frac{4L}{\lambda n}
\label{eq:three}
\end{align}
If we perturb the data by a single element, we learn $\mathcal{A}$ that can become arbitrarily close for large $n$.

We then use \ref{eq:three} and the $L$-Lipschitz property of $l(\cdot,z)$:
$$\implies \sup_z[l(\mathcal{A}(S), z) - l(\mathcal{A}(S_{i,z_i'}), z)| \leq \frac{4L^2}{\lambda n}$$
\end{proof}


\section*{Stochastic Gradient Descent (SGD) is Stable}

\subsection*{Stability Theorem}
Recall the SGD update formula,

\begin{equation}
    w_{t+1} = w_t - \alpha_t \nabla_w l(w_t, z_{i,t}), i_t   \sim \text{uniform}(1, \cdots, n )
\end{equation}

where $w_t$ is the weight iterate at time $t$, $\alpha_t$ is an (annealing) learning rate at time $t$ and $l(w_t, z_{i,t})$ is the computed loss for the current weight iterate for a particular example $z_{i,t}$.
 

\begin{thm}
If $f(\cdot, z)$ is $\gamma$-smooth, convex and $L$-Lipschitz, then Stochastic Gradient Descent is $\beta$-uniformly stable where

\[
\beta \leq \frac{2L^2}{n} \sum_{t=1}^{T} \alpha_t
\]
\end{thm}

\textbf{Analysis:} \\
We are no longer requiring the function to be strongly convex.  Additionally, this result holds for a finite number of steps $T$.  

\subsection*{Stability Proof (Rough Outline)}
We will consider two runs of the SGD algorithm.  One run will be on the original data set $S$ and the other run will be on the data set $S_{i,z_i'}$.  Recall, this indicates the same data set $S$ only now with the $i^{th}$ element swapped with element $z_i'$.  In order to compare the stability between the two runs, we maintain the same order of element selection (same random seed) for $t=1, \cdots, T$.

\vspace{0.2cm}

\begin{defn} 
    \[
    \delta_t = || w_t - w_t'||
    \]
\end{defn}


where $w_t'$ denotes the iterate for the SGD algorithm on the data set $S_{i,z_i'}$.

We can write the expectation of the difference $\delta_{t+1}$ as the following:

\begin{equation} \label{eqn: sgd_prob_decomp}
    E[\delta_{t+1}] = P(i_t = i) E[\delta_{t+1} | i_t = i] + P(i_t \neq i) E[\delta_{t+1} | i_t \neq i]
\end{equation}

We introduce two Lemmas
\begin{lemma} \label{lem: lem_1}
    \[
    E[\delta_{t+1} | i_t \neq i] \leq E[\delta_t]
    \]
\end{lemma}
\begin{proof}
Convexity and $\gamma$-smoothness implies that the gradients are co-coercive for a  function $f$: 
\[ 
\langle\,\nabla f(v)- \nabla f(w), v-w \rangle \geq \frac{1}{\gamma}||\nabla f(v)- \nabla f(w)||^2
\]
We conclude that the weight update can be expressed as:
    \[ 
    ||w_{t+1}-w'_{t+1}||^2 = ||w_t-w'_t||^2 - 2\alpha_t \langle \nabla f(w_t)- \nabla f(w'_t), w_t-w'_t \rangle +\alpha^2 ||\nabla f(w_t)- \nabla f(w'_t)||^2 \]
    \[
    \leq ||w_t-w'_t||^2 - (2 \alpha_t / \gamma - \alpha_t^2) ||\nabla f(w_t)- \nabla f(w'_t)||^2 \leq ||w_t-w'_t||^2
    \]
so we get, using definition 7 that:
    \[
    ||w_{t+1}-w'_{t+1}|| = \delta_{t+1} \leq ||w_{t}-w'_{t}|| = \delta_{t}
    \]
\end{proof}
\begin{lemma} \label{lem: lem_2} And for the index that has been swapped
    \[
    E[\delta_{t+1} | i_t = i] \leq E[\delta_t] + 2 \alpha_t L
    \]
where $L$ is the Lipschitz value.
\end{lemma}
\begin{proof}
We know that 
    \[\delta_{t+1}=||w_{t+1}-w_{t+1}'||=||w_t-\alpha_t \nabla l(w_t, z_{i_t}) - ( w_t'-\alpha_t \nabla l(w_t', z_{i_t}))||
    \]
Using the triangle inequality we can write
    \[\delta_{t+1} \leq ||w_t - w_t'|| + \alpha_t||\nabla l(w_t,z_{i_t}) - \nabla l(w_t',z_{i_t})||
    \]
Since $l(\cdot, z)$ is L-lipschitz
    \[\delta_{t+1} \leq \delta_t + 2\alpha_t L
    \]
Taking expectation on either side we get
    \[E[\delta_{t+1}|i_t=i] \leq E[\delta_t] + 2\alpha_t L
    \]
\end{proof}

  Using Lemmas \ref{lem: lem_1}, \ref{lem: lem_2}, we may rewrite Equation \ref{eqn: sgd_prob_decomp} as:

\begin{equation}
    E[\delta_{t+1}] \leq \left(1 - \frac{1}{n}\right)E[\delta_t] + \frac{1}{n} \left( E[\delta_t] + 2 \alpha_t L \right)
\end{equation}

which when recursively unrolled yields the following final $\delta_T$

\begin{equation}
    E[\delta_T] = E[|| w_T - w_T' ||] \leq  \sum_{t=0} ^ {T-1}  \frac{ 2 \alpha_t L }{n}
\end{equation}
We find that:
    \[
    E[\delta_{t+1}] = P(i_t=i)E[\delta_{t+1}|i_t=i] + P(i_t \neq i)E[\delta_{t+1}|i_t \neq i]] 
    \]
    \[
    \leq \frac{1}{n}(E[\delta_t] + 2\alpha_t L) + E[\delta_t](1+\frac{1}{n}) \leq E[\delta_t] + \frac{2 \alpha_t L}{n}
    \]
SGD is therefore \textbf{stable} since $\sum_{t=0} ^ {T-1}  \frac{ 2 \alpha_t L }{n} \equiv \beta$ is $O(\frac{1}{n})$ for $n$ data points.



%----------------------------------------

%----------------------------------------
\bibliographystyle{abbrvnat}
\bibliography{../Refs/lec9b}
%----------------------------------------
\end{document}
